\chapter{Обзор бесступенчатых коробок передач}

Бесступенчатые коробки передач (вариаторы, CVT — Continuously Variable Transmission) представляют собой трансмиссии, в которых передаточное отношение изменяется плавно, без фиксированных ступеней. Это позволяет двигателю постоянно работать в оптимальном диапазоне оборотов, что повышает топливную экономичность и снижает динамические нагрузки на узлы трансмиссии. В современной практике применяются несколько основных разновидностей вариаторов, отличающихся по принципу передачи крутящего момента.

\section{Клиновый вариатор}

Клиновый вариатор (рис. \ref{img:Vcvt}) относится к типу бесступенчатых трансмиссий, в которых передача крутящего момента осуществляется за счёт трения между парой раздвижных шкивов и гибким элементом клинового профиля — ремнём или цепью. Изменение расстояния между половинами шкивов плавно изменяет их рабочие диаметры, обеспечивая непрерывное регулирование передаточного числа без рывков и переключений передач.

\includeimage
    {Vcvt} % Имя файла без расширения (файл должен быть расположен в директории inc/img/)
    {f} % Обтекание (без обтекания)
    {h} % Положение рисунка (см. figure из пакета float)
    {0.5\textwidth} % Ширина рисунка
    {Конуса клиновидного типа.} % Подпись рисунка

Такие вариаторы отличаются плавностью, компактностью и высокой эффективностью, особенно при малых и средних нагрузках. В зависимости от конструкции гибкого элемента различают три основных типа клиновых вариаторов — с сухим ремнём (Dry Belt), металлической цепью (Chain) и толкающим ремнём (Pushbelt). Каждый из них имеет собственные особенности, определяющие сферу применения — от лёгких мотоциклов до современных автомобилей повышенной мощности.

\subsection{Сухой ремень}

Развитие вариаторов клинового типа началось с применения резиновых клиновых ремней. Такие ремни и сегодня находят применение, однако их ограниченная способность передавать высокий крутящий момент делает их малопригодными для автомобильных двигателей большой мощности.

Особый интерес представляют сухие ремни (рис. \ref{img:DryBelt}), в которых отсутствует смазка. Благодаря этому между ремнём и шкивами устанавливается повышенный коэффициент трения, что позволяет уменьшить прижимное усилие, сделать вариатор компактнее и легче. Такие конструкции хорошо подходят для легких мотоциклов и малолитражных автомобилей, где требуемая мощность невелика.

\includeimage
    {DryBelt} % Имя файла без расширения (файл должен быть расположен в директории inc/img/)
    {f} % Обтекание (без обтекания)
    {h} % Положение рисунка (см. figure из пакета float)
    {0.5\textwidth} % Ширина рисунка
    {Сухой ремень.} % Подпись рисунка

Главный недостаток сухих ремней — отсутствие охлаждения зоны контакта из-за отсутствия масла. Это приводит к перегреву и ограничивает передаваемый крутящий момент, снижая долговечность при интенсивной эксплуатации.

\subsection{Цепь}

Цепные клиновые вариаторы (рис. \ref{img:Chain}) состоят из сегментов и штифтов, образующих гибкую металлическую цепь. Штифты имеют сферическую поверхность, что улучшает контакт со шкивами и снижает износ. Внутри цепи предусмотрены вращательные контакты, исключающие проскальзывание, что обеспечивает высокий КПД, особенно при низких крутящих моментах.

\includeimage
    {Chain} % Имя файла без расширения (файл должен быть расположен в директории inc/img/)
    {f} % Обтекание (без обтекания)
    {h} % Положение рисунка (см. figure из пакета float)
    {0.5\textwidth} % Ширина рисунка
    {Цепной ремень.} % Подпись рисунка

У цепных вариаторов почти отсутствует внутреннее трение, поэтому они работают эффективнее, чем толкающие ремни, хотя и шумнее. Повышенный уровень шума связан с малым числом штифтов, которые поочередно входят в зацепление со шкивом.

Цепные вариаторы обладают высокой прочностью и долговечностью, но при этом более сложны в производстве, тяжелее и дороже ременных аналогов.

\subsection{Толкающий ремень}

Толкающий ремень (рис. \ref{img:PushBelt}) состоит из металлических блоков (рис. \ref{img:PushBeltPlate}), удерживаемых несколькими пакетами тонких стальных лент (обычно 9–12 в каждом пакете). Ленты работают на растяжение, а блоки — на сжатие, благодаря чему ремень способен передавать крутящий момент в режиме сжатия, что и дало название системе.

\includeimage
    {PushBelt} % Имя файла без расширения (файл должен быть расположен в директории inc/img/)
    {f} % Обтекание (без обтекания)
    {h} % Положение рисунка (см. figure из пакета float)
    {0.5\textwidth} % Ширина рисунка
    {Толкающий ремень.} % Подпись рисунка

Такая конструкция обеспечивает высокую жёсткость и плавность хода, а также более тихую работу по сравнению с цепными вариаторами, поскольку количество блоков значительно больше числа звеньев цепи. Это снижает давление между блоками и шкивами и увеличивает износостойкость.

\includeimage
    {PushBeltPlate} % Имя файла без расширения (файл должен быть расположен в директории inc/img/)
    {f} % Обтекание (без обтекания)
    {h} % Положение рисунка (см. figure из пакета float)
    {0.5\textwidth} % Ширина рисунка
    {Блоки и ленты толкающего ремня.} % Подпись рисунка

Однако при работе толкающий ремень испытывает значительные потери на трение, вызванные разницей скоростей между блоками и лентами, а также усталостные нагрузки из-за постоянных изгибов и растяжений. Эти факторы ограничивают максимальный передаваемый момент и требуют жёстких требований к материалам и смазке. Несмотря на это, pushbelt остаётся наиболее совершенной и сбалансированной конструкцией клинового вариатора, применяемой в современных автомобилях среднего и премиум-класса.

\subsection{Обзор применения клиновых вариаторов}

На представленном графике (рис. \ref{img:CVTBeltMarket}) распределения долей рынка по типам ремней для вариаторов видно, что толкающий ремень (push belt) занимает около 48 \%, в то время как цепной тип (chain CVT) — примерно 52 \%. Такое соотношение говорит о почти равной востребованности двух решений, однако с некоторым преимуществом цепных систем, используемых преимущественно в трансмиссиях автомобилей среднего и премиального сегмента. Толкающие ремни чаще применяются в массовых моделях, где требуется баланс между стоимостью, надёжностью и эффективностью.

\includeimage
    {CVTBeltMarket} % Имя файла без расширения (файл должен быть расположен в директории inc/img/)
    {f} % Обтекание (без обтекания)
    {h} % Положение рисунка (см. figure из пакета float)
    {0.5\textwidth} % Ширина рисунка
    {Анализ рынка типов ремней для клиновидных вариаторов \cite{CVTBeltMarketLib}} % Подпись рисунка

\section{Тороидальный вариатор}

Тороидальный вариатор (рис. \ref{img:Ocvt}) передаёт крутящий момент через ролики, расположенные между тороидальными дисками. Изменение угла наклона роликов изменяет эффективный радиус контакта и, соответственно, передаточное число. Преимущества включают высокую передаваемую мощность, равномерное распределение нагрузки и долговечность при интенсивной эксплуатации. Недостатки — сложность конструкции, высокая стоимость изготовления и необходимость точной смазки. Такие вариаторы применяются в премиальных автомобилях и гибридных системах, где требуется сочетание высокой мощности и плавного изменения передаточного отношения.

\includeimage
    {Ocvt} % Имя файла без расширения (файл должен быть расположен в директории inc/img/)
    {f} % Обтекание (без обтекания)
    {h} % Положение рисунка (см. figure из пакета float)
    {0.5\textwidth} % Ширина рисунка
    {Диски и ролики тороидального вариатора.} % Подпись рисунка

\section{Конусный вариатор}

Конусный вариатор (рис. \ref{img:Wcvt}) передаёт крутящий момент через ролик или ремень, контактирующий с двумя конусными поверхностями. Передаточное отношение регулируется смещением точки контакта вдоль оси конусов. Преимущества — простота конструкции, компактность и высокая плавность регулирования. Однако он рассчитан на сравнительно малые нагрузки, имеет повышенный износ контактных поверхностей и чувствителен к загрязнению и смазке. Эти ограничения делают конусные вариаторы непригодными для большинства легковых автомобилей, где требуется долговечность и возможность передавать значительные крутящие моменты. Конусные вариаторы находят применение в маломощных транспортных средствах, скутерах, мотоциклах и промышленном оборудовании.

\includeimage
    {Wcvt} % Имя файла без расширения (файл должен быть расположен в директории inc/img/)
    {f} % Обтекание (без обтекания)
    {h} % Положение рисунка (см. figure из пакета float)
    {0.5\textwidth} % Ширина рисунка
    {Конуса и ремень конусного вариатора.} % Подпись рисунка

\section{Электрический вариатор}

Электрический вариатор (рис. \ref{img:eCVT}) используется в гибридных автомобилях и сочетает двигатель внутреннего сгорания с двумя электродвигателями и планетарным редуктором. ДВС подключается к солнечной шестерне, один электродвигатель — к кольцевой шестерне, а второй — к планетарным роликам или вала выхода на колёса. Такая схема позволяет распределять мощность между колесами и электродвигателями и плавно изменять передаточное число без дискретных передач. Передаточное число регулируется за счёт изменения соотношения вращения ДВС и электродвигателей, что удерживает двигатель в оптимальной зоне оборотов. Система обеспечивает плавность хода, мгновенный отклик на изменение нагрузки и рекуперацию энергии при торможении, повышая экономичность и комфорт движения.

\includeimage
    {eCVT} % Имя файла без расширения (файл должен быть расположен в директории inc/img/)
    {f} % Обтекание (без обтекания)
    {h} % Положение рисунка (см. figure из пакета float)
    {0.5\textwidth} % Ширина рисунка
    {Генератор (слева), планетарный редуктор и электодвигатель (справа) электрического вариатора.} % Подпись рисунка

\section{Гидростатический вариатор}

Гидростатический вариатор передаёт крутящий момент с помощью насоса и гидромотора с регулируемым рабочим объёмом. Изменение подачи рабочей жидкости позволяет плавно изменять передаточное число. Преимущества — способность передавать значительные нагрузки, надёжность и возможность точной регулировки скорости. Недостатки — меньший КПД по сравнению с механическими системами и сложность конструкции. Гидростатические вариаторы находят применение в сельскохозяйственной технике, промышленном оборудовании и специализированных транспортных средствах, где требуется высокая точность управления движением.

Таким образом, бесступенчатые трансмиссии представлены несколькими конструктивными схемами, каждая из которых имеет свои преимущества и ограничения. Выбор типа вариатора определяется назначением транспортного средства, требуемой мощностью, условиями эксплуатации и требованиями к эффективности.

\section{Анализ рынка бесступенчатых коробок передач}

Анализ рынка бесступенчатых коробок передач (рис. \ref{img:CVTMarket}) демонстрирует, что клинноременные вариаторы занимают доминирующее положение в сегменте легковых автомобилей благодаря своей эффективности и экономичности. Основные производители, такие как Jatco, Aisin и ZF, активно развивают технологии клиноремённых вариаторов, внедряя инновации для повышения долговечности и снижения износа ремней.

\includeimage
    {CVTMarket} % Имя файла без расширения (файл должен быть расположен в директории inc/img/)
    {f} % Обтекание (без обтекания)
    {h} % Положение рисунка (см. figure из пакета float)
    {0.5\textwidth} % Ширина рисунка
    {Анализ рынка автомобилей с бесступенчатой трансмиссией \cite{CVTMarketLib}} % Подпись рисунка

Как показывают исследования в рамках текущего курсового проекта стоит выбрать клиновый вариатор с толкающим ремнём. Эта комбинация обеспечивает оптимальное соотношение между надежностью, долговечностью и экономичностью. Толкающий ремень способен передавать относительно высокие крутящие моменты при меньших потерях, отличается плавностью работы и сравнительно низким уровнем шума. Клиноременная конструкция вариатора проста в производстве и эксплуатации, что делает её предпочтительной для массового применения в автомобилях среднего сегмента.